\label{sec:legal}

\subsection{Communities of Interest in Redistricting Law}

Economic community membership is one of the oldest recognized forms of
``communities of interest'' in redistricting law.
The doctrine requires that district lines not unnecessarily divide
communities whose members share common social, cultural, economic, or
geographic interests \citep{ncsl2021criteria}.

Several states have codified economic communities as an explicit criterion:

\paragraph{California Proposition 11 (2008).}
The Voters First Act, which established California's independent Citizens
Redistricting Commission, lists ``communities of interest'' as a primary
criterion and explicitly identifies economic interests as a component
\citep{caprop11}.
The Commission's implementation guidelines have treated shared commercial
corridors, agricultural communities, and employment districts as qualifying
economic communities.

\paragraph{Colorado Constitution, Article V, Section 44.}
Colorado's independent redistricting commission criteria require that
districts preserve communities of interest to the extent possible, with
guidance that explicitly includes economic communities \citep{coart5s44}.

\paragraph{Arizona Proposition 106 (2000).}
Arizona's Independent Redistricting Commission criteria list
``communities of interest'' including economic and social interests \citep{azircprop106}.

\paragraph{Washington RCW 44.05.090.}
Washington's redistricting law requires preservation of ``communities of
related and mutual interests,'' with economic community membership
recognized in Commission practice \citep{warcw44}.

At the federal level, \textit{Rucho v.\ Common Cause} \citeyearpar{rucho2019}
established that federal courts cannot adjudicate partisan gerrymandering claims.
Economic character weights are relevant to this context because they provide
a \emph{partisan-neutral} rationale for district lines: the weight
modification uses no electoral data, no voter registration data, and no
racial or ethnic composition data.
The inputs are NAICS-sector job counts from a federal government
administrative dataset (LODES WAC) that are published annually by the
Census Bureau's LEHD program.

\subsection{The Partisan-Neutrality Defense}

The legal defensibility of algorithmic redistricting depends in part on
demonstrating that the inputs to the algorithm are partisan-neutral.
Economic character weights satisfy this requirement structurally:

\begin{itemize}
  \item \textbf{Data source.} LEHD LODES WAC data reports workplace
    employment counts by NAICS sector.
    No racial, ethnic, or partisan data is present in the WAC file.
  \item \textbf{Computation.} The economic character vector
    $\mathbf{e}(t) = (\mathrm{CI}(t), \mathrm{IF}(t), \mathrm{JPR}(t))$
    is computed purely from job counts and population.
    No electoral geography, no voter file, no precinct boundary data enters
    the computation.
  \item \textbf{Causation direction.} The economic character signal
    influences district boundaries only through the economic profile of
    tracts, not through their partisan voting history.
    Any partisan effect is a consequence of the geographic correlation
    between economic activity and partisan sorting---a correlation that
    exists independently of the algorithm's choices.
\end{itemize}

This structure allows an expert witness to testify that ``district boundaries
follow natural economic zone transitions, not partisan boundaries,'' which is
legally stronger than asserting a proportionality outcome.

\subsection{Expert Witness Framing}

For expert witnesses using economic character weights in a redistricting
proceeding, we recommend the following framing elements:

\begin{enumerate}
  \item \textbf{Data provenance.} ``The edge weight modification uses the
    Census Bureau's LODES Workplace Area Characteristics data, a federal
    administrative dataset published annually by the LEHD program.
    The data records the number of jobs by NAICS sector in each census block.
    No electoral, racial, or demographic data enters the weight computation.''

  \item \textbf{Criterion alignment.} ``Preserving economic community character
    serves the [state]'s stated redistricting criterion of maintaining
    communities of interest.
    Tracts with similar economic profiles---both residential, both commercial,
    both industrial---are kept together where possible; district cuts are
    drawn at economic zone transitions, which are natural community boundaries.''

  \item \textbf{Mechanistic transparency.} ``The weight formula assigns
    a cosine similarity score between 0 and 1 to each pair of adjacent tracts
    based on their economic character vectors.
    High similarity (similar economic profile) makes the edge harder to cut;
    low similarity (different economic profiles) makes the edge easier to cut.
    The $\alpha = 0.5$ blend factor ensures that geographic contiguity
    remains the dominant consideration.''

  \item \textbf{State-specificity.} ``In North Carolina, this modification
    reduces the degree to which major employment centers (Research Triangle
    Park, Charlotte commercial corridor) are embedded within residential
    partisan districts, producing districts that more cleanly separate
    employment zones from residential zones.
    In states without concentrated employment clusters, the modification
    is appropriately neutral.''
\end{enumerate}

\subsection{When to Use Economic Character Weights}

Economic character weights are most valuable, and most legally supportable,
in states that satisfy the following conditions:

\begin{itemize}
  \item The state has an explicit communities-of-interest criterion in its
    redistricting law or commission criteria.
  \item The state has concentrated employment centers (university research
    parks, hospital campuses, industrial corridors) that are currently being
    split by district lines drawn primarily on geographic compactness grounds.
  \item LODES data coverage for the state is complete (all states are covered).
  \item The practitioner can demonstrate, from the LODES data, that the
    employment concentration pattern exists and that the proposed plan
    respects it.
\end{itemize}

Conversely, in states like Wisconsin where economic geography is diffuse,
economic character weights add no value and should not be used as a
justification for a specific plan; the null result will undermine rather
than support the communities-of-interest argument.
